\documentclass[conference,a4paper]{IEEEtran}
\usepackage{cite}
\usepackage{amsmath,amssymb,amsfonts}
\usepackage{graphicx}
\usepackage{textcomp}
\usepackage{xcolor}
\usepackage{hyperref}
\usepackage{lettrine}

\bibliographystyle{IEEEtran}

\IEEEoverridecommandlockouts
\begin{document}

\title{Backside Power Delivery for Sub-5 nm CMOS: Signal Integrity Benefits and Emerging Design Trade-offs}

\author{\IEEEauthorblockN{Daniel Tyukov (5714699)\thanks{This article was written as part of the course EE4C01 Profile Orientation \& Academic Skills of the faculty Electrical Engineering, Mathematics and Computer Science of Delft University of Technology. Submission date: April 9, 2026. Teacher: Matthijs Looij. Total word count: 2000.}}
\IEEEauthorblockA{Department of Microelectronics\\
Delft University of Technology\\
Delft, The Netherlands\\
d.tyukov@student.tudelft.nl}
\and
\IEEEauthorblockN{Apostolos Tsailas (5506387)}
\IEEEauthorblockA{Department of Microelectronics\\
Delft University of Technology\\
Delft, The Netherlands\\
a.tsailas-1@student.tudelft.nl}
}

\maketitle

\begin{abstract}
Technology scaling beyond 5\,nm has revealed fundamental limitations in conventional front-side power delivery networks (FSPDNs). As interconnect dimensions shrink, increased wire resistance and routing congestion make it progressively harder to maintain power integrity without sacrificing signal routing resources and overall performance. These constraints have motivated the development of back-side power delivery networks (BSPDNs), which relocate the power grid to the wafer backside and decouple it from the signal interconnect entirely.

This literature review examines the transition from FSPDNs through buried power rails to full backside power delivery in sub-5\,nm CMOS technologies. Nine peer-reviewed sources, published between 2019 and 2025 and drawn from IEEE Xplore, were analysed and compared. The review evaluates reported improvements in IR~drop, voltage droop, operating frequency, and cell area, and contrasts these gains against new challenges introduced by the BSPDN architecture.

The reviewed studies demonstrate that BSPDNs reduce peak IR~drop by more than four times compared to conventional FSPDNs, improve voltage droop by up to 30\%, and enable cell area savings of approximately 19\%. However, backside delivery introduces higher thermal resistance, with hotspot temperatures increasing by up to 45\%, and raises concerns around self-heating and AC impedance resonance. Overall, BSPDNs represent a promising architectural evolution for power delivery at advanced technology nodes, provided that thermal and fabrication challenges are addressed through co-design strategies.
\end{abstract}

\section{Introduction}

\lettrine{T}{echnology} scaling beyond 5nm has introduced significant challenges in power delivery and interconnect performance. Although transistor dimensions continue to shrink, the scaling of interconnects has not kept pace, leading to increased wire resistance and tighter voltage margins. Conventional front-side power delivery networks (PDNs), which share routing resources with signal interconnects, face growing difficulty in meeting power integrity requirements without degrading routing efficiency and overall performance\cite{8993617}.

To address such limitations buried power rails (BPR) have been proposed as an early innovation. By relocating power rails beneath active devices and utilizing high-aspect-ratio ruthenium structures, BPR reduces local resistance and frees front-side routing tracks, enabling improved IR-drop characteristics\cite{8993617}. Backside Power Delivery Networks (BS-PDN) move the power grid to the wafer backside which decouples the power and signal routing\cite{10185208}. Experimental findings have shown substantial reductions in voltage droop and measurable frequency improvements in high-performance cells \cite{10185208}. Additionally, significant area savings and enhanced routing efficiency have been demonstrated in both high-performance and high-density scenarios\cite{10631463}, while models indicate more than a four-fold reduction in IR~drop compared to conventional PDNs\cite{8932458}.

These architectural innovations represent a fundamental shift in power delivery design, moving from incremental BEOL optimizations toward structural reorganization of the power distribution network. This literature review therefore examines the evolution from conventional front-side PDNs to buried power rails and backside power delivery, evaluating their signal integrity and power--performance--area benefits as well as the emerging design challenges they introduce in sub-5\,nm CMOS technologies. The review first outlines the limitations of conventional front-side PDNs, then presents the improvements achieved by backside delivery, followed by an analysis of the new challenges this architecture introduces, and concludes with a discussion of open questions and directions for future research.\section{Introduction}

\lettrine{T}{echnology} scaling beyond 5nm has introduced significant challenges in power delivery and interconnect performance. Although transistor dimensions continue to shrink, the scaling of interconnects has not kept pace, leading to increased wire resistance and tighter voltage margins. Conventional front-side power delivery networks (PDNs), which share routing resources with signal interconnects, face growing difficulty in meeting power integrity requirements without degrading routing efficiency and overall performance\cite{8993617}.

To address such limitations buried power rails (BPR) have been proposed as an early innovation. By relocating power rails beneath active devices and utilizing high-aspect-ratio ruthenium structures, BPR reduces local resistance and frees front-side routing tracks, enabling improved IR-drop characteristics\cite{8993617}. Backside Power Delivery Networks (BS-PDN) move the power grid to the wafer backside which decouples the power and signal routing\cite{10185208}. Experimental findings have shown substantial reductions in voltage droop and measurable frequency improvements in high-performance cells \cite{10185208}. Additionally, significant area savings and enhanced routing efficiency have been demonstrated in both high-performance and high-density scenarios\cite{10631463}, while models indicate more than a four-fold reduction in IR~drop compared to conventional PDNs\cite{8932458}.

These architectural innovations represent a fundamental shift in power delivery design, moving from incremental BEOL optimizations toward structural reorganization of the power distribution network. This literature review therefore examines the evolution from conventional front-side PDNs to buried power rails and backside power delivery, evaluating their signal integrity and power--performance--area benefits as well as the emerging design challenges they introduce in sub-5\,nm CMOS technologies. The review first outlines the limitations of conventional front-side PDNs, then presents the improvements achieved by backside delivery, followed by an analysis of the new challenges this architecture introduces, and concludes with a discussion of open questions and directions for future research.
\section{Methodology}

Literature was collected from IEEE Xplore. Search terms combined ``backside power delivery,'' ``buried power rails,'' ``front-side PDN,'' ``IR-drop,'' ``signal integrity,'' and ``sub-5nm CMOS.'' More specific terms such as ``nanoTSV,'' ``PowerVia,'' and ``power-performance-area'' were added to capture implementation studies. Only publications from 2019 onward were considered, since the earliest demonstrations of buried power rails and backside delivery at advanced nodes date from that year.

The initial search returned over 80 results. Titles and abstracts were screened first: papers had to deal directly with power delivery network architecture and its effect on power integrity or routing efficiency at sub-5\,nm nodes. Work focused purely on transistor-level scaling or on unrelated packaging topics was excluded. The remaining papers were then read in full. Each was judged on whether its results rested on simulation, measurement, or both, and on how much it added to the comparison between front-side and backside delivery. Papers that reported quantitative IR-drop, voltage droop, or area figures across different PDN configurations were preferred.

Nine sources were retained. They span early architecture proposals, industry demonstrations from Intel and imec, and parasitic modelling of backside configurations. Together they cover the path from conventional front-side PDN through buried power rails to full backside delivery, with both measured and simulated data on the trade-offs involved.

\section{Limitations of Conventional Front-Side Power Delivery Networks}

Conventional integrated circuits rely on front-side power delivery (FSPDNs) where power and signal interconnects are routed via the same back-end-of-line (BEOL) metal stack. In this layout, supply voltages are distributed from the package connections through all individual metal layers, down to the transistor level. This architecture has been used for multiple generations of integrated circuits, however continued scaling below 5 nm has revealed several limitations that have made reliable power delivery more difficult \cite{9817394},\cite{10244504}. 

As a result of the scaling of interconnects, an issue arises from the increased resistance. As metal dimensions shrink, resistivity increases due to surface and grain-boundary scattering, leading to larger IR~drops across the power grid\cite{8993617}. At the local metal layers closest to the transistors, wire widths fall below 20\,nm, and at these dimensions copper resistivity rises sharply because the wire is narrower than the electron mean free path\cite{8932458}. This results in a reduced effective supply voltage at the transistors, which can degrade circuit performance and reliability. In the most advanced node, highly resistive local metal layers contribute significantly to IR-drop hotspots, particularly in areas with dense logic \cite{8993617}. Studies conducted with an ARM Cortex-A53 design at a 3\,nm technology node show that even with an optimised power grid, conventional FSPDNs did not meet the 10\% V\textsubscript{DD} IR-drop target\cite{9817394}.

Another key limitation that arises due to scaling is the increased competition for the limited routing resources available for power and signal interconnects. In FSPDNs, both power rails and signal interconnects occupy the same metal stack. As such, designers must dedicate a large portion of routing resources to maintain a sufficiently reliable power grid that is able to provide the necessary performance. This results in reduced routing resources for signal interconnects and increased routing congestion. In high-performance designs, the power grid can take up four or more metal layers that would otherwise carry signals, which forces designers to choose between a larger die area or lower performance\cite{10631463}. Block-level studies at imec showed that high-performance standard-cell libraries need an 18-CPP power stripe pitch on the front side just to keep IR~drop acceptable, while backside delivery reaches the same power integrity without using any front-side routing tracks at all\cite{10631463}.

To mitigate these limitations, Buried Power Rails (BPR) are used as a current solution. These designs move power rails beneath active devices in order to reduce resistance and make more routing resources available to the signal interconnect. BPR-based designs have shown a 25\% reduction in on-chip IR drop and 17\% reduction in off-chip voltage droop compared with conventional front-side PDNs \cite{9817394}. However, even with these improvements, standalone BPR designs are viewed as intermediate solutions\cite{10244504}. The fundamental limitations of front-side routing remain.

Backside power delivery networks (BSPDNs) have been explored as a long-term solution to these issues. In this architecture, the entire power grid is moved to the backside of the wafer, typically through nano through-silicon vias (nTSVs) that connect the backside metal layers directly to the transistor contacts. By fully separating power and signal routing onto opposite sides of the die, BSPDN architecture eliminates the resource conflict between the two and enables both lower IR~drop and improved routing density simultaneously. Multiple studies have confirmed the viability of this approach at 3\,nm and below\cite{10631463},\cite{8932458},\cite{10244504}.


\section{Improvements of Back-Side Power Delivery Network Implementations}

Back-Side Power Delivery Networks improve power integrity and routing efficiency by relocating the power grid to the backside of the wafer and utilising the front-side solely for signal interconnect. The design allows for a reduction in congestion and increases routing flexibility, especially in advanced technology nodes\cite{10631463}. In dense logic areas, this architecture allows for signal interconnects to utilize lower layers on the front side of the wafer and, as a result, reduces wire lengths, reduces interconnect delay, and improves timing\cite{11192533}.

A primary improvement of the BSPDN architecture is the reduction in IR drop. By delivering power from a dedicated backside metal layer, voltage droop at the transistor level is reduced. Modelling of backside PDN configurations incorporating buried power rails and µTSVs demonstrated a more than four-fold reduction in IR-drop compared to conventional front-side PDNs\cite{8932458}. This improvement is due to the lower resistance of the thicker backside power rails and the shorter current path to the transistor layer, which bypasses the highly resistive local BEOL metal layers entirely.

Intel's PowerVia Technology, demonstrated using a fabricated E-core, showed that a BSPDN defforsign yielded a 30\% improvement in voltage droop and 6\% frequency improvement compared to a reference FSPDN design\cite{10185208}. By splitting power delivery from the signal interconnect, PowerVia supplies a more uniform supply voltage across the die, resulting in lower average and worst-case IR-drop hotspots that are particularly problematic in dense logic regions\cite{10185208}.

BSPDN architecture has also shown improvements in routing efficiency and cell area. With the power grid moved to the backside of the wafer, front-side routing tracks that would be used for power delivery are freed, enabling either area reduction or increased routing flexibility at the same cell dimensions. Studies evaluating high-density designs showed area savings up to 19\% when implementing a BSPDN architecture. Signal interconnect was completed at lower metal layers, reducing overall wire length\cite{10631463}.

Table~\ref{tab:comparison} summarises the quantitative improvements reported across the reviewed studies for each power delivery architecture relative to the conventional FSPDN baseline.

\begin{table}[h]
\centering
\caption{Comparison of power delivery architectures at sub-5\,nm nodes. Values indicate improvement relative to conventional FSPDN.}
\label{tab:comparison}
\begin{tabular}{|l|c|c|c|}
\hline
\textbf{Metric} & \textbf{FSPDN} & \textbf{BPR} & \textbf{BSPDN} \\
\hline
On-chip IR-drop & Baseline & $-$25\%\cite{9817394} & $>$4$\times$ reduction\cite{8932458} \\
\hline
Off-chip voltage droop & Baseline & $-$17\%\cite{9817394} & $-$30\%\cite{10185208} \\
\hline
Frequency & Baseline & --- & +6\%\cite{10185208} \\
\hline
Cell area & Baseline & --- & $-$19\%\cite{10631463} \\
\hline
Thermal (hotspot) & Baseline & --- & +45\%\cite{10925422} \\
\hline
\end{tabular}
\end{table}
\section{Limitations of Back-Side Power Delivery Networks}

The improvements in power integrity and routing efficiency offered by BSPDNs come at a cost. Moving the power grid to the backside of the wafer introduces thermal, self-heating, and AC impedance challenges that are not present in conventional front-side architectures.

The most pressing issue is thermal performance. In FSPDNs, the full-thickness silicon substrate spreads heat laterally from the active devices toward the topside heat sink. BSPDNs require thinning the wafer to form nano through-silicon via (nTSV) connections, which removes much of this lateral heat-spreading path. The bonding interface and back-end-of-line (BEOL) layers between the devices and the heat sink further increase thermal resistance. Modelling of CPU hotspot regions has shown that BSPDN configurations reach temperatures approximately 45\% higher than equivalent FSPDN designs\cite{10925422}. Several mitigation strategies have been investigated, including high thermal conductivity bonding materials, stacked BEOL via configurations that create direct vertical thermal paths, and embedded microchannel cooling etched into the backside metal layers\cite{10925422}.

At the device level, the thinned substrate also worsens self-heating effects (SHE). With less silicon to absorb and spread the heat generated by switching transistors, localised temperature rises become more severe. An evaluation of SHE in SRAM macros and circuits at the 3\,nm node found that BSPDN designs improve power, performance, and area, but the resulting thermal behaviour requires power-thermal co-design to avoid reliability degradation\cite{10529350}. The study further showed that the choice of buried power rail material and the number of backside metal layers both have a measurable impact on device temperature, indicating that thermal optimisation must be considered alongside electrical design choices from the earliest stages of the design flow\cite{10529350}. Electrical gains therefore cannot be evaluated in isolation from their thermal consequences.

A separate concern relates to the AC impedance profile of the backside power grid. Because the DC resistance of BSPDNs is much lower than that of FSPDNs, the damping in the power delivery network is reduced. As a result, impedance resonance peaks can shift into frequency ranges that coincide with circuit switching activity, which risks supply voltage fluctuations during transient load changes\cite{8993617}. Studies on the Arm Cortex-A53 benchmark showed that this resonance effect is particularly relevant in multi-core designs, where on-chip decoupling capacitors are needed to suppress ringing\cite{8993617}. Careful impedance profiling and decoupling capacitor placement during the design phase are therefore essential when adopting a backside power architecture.


\section{Discussion and Conclusion}

This literature review has examined the evolution of power delivery architectures for sub-5\,nm CMOS, from conventional front-side PDNs through buried power rails to full backside delivery. The reviewed studies consistently demonstrate that BSPDNs provide substantial improvements in power integrity and design efficiency. By relocating the power network to the wafer backside, this architecture decouples power and signal routing, enabling IR~drop reductions of more than four times\cite{8932458}, voltage droop improvements of 30\%\cite{10185208}, frequency gains of 6\%\cite{10185208}, and cell area savings of up to 19\%\cite{10631463}. These results, drawn from both simulation and fabricated test vehicles, indicate that backside delivery addresses the fundamental scaling limitations that front-side PDNs face at advanced nodes. Buried power rails serve as an effective intermediate step, reducing IR~drop by 25\% and voltage droop by 17\%\cite{9817394}, but the full separation of power and signal routing achieved by BSPDNs delivers substantially greater benefits across all evaluated metrics.

However, these gains come with trade-offs that cannot be overlooked. The thinned substrate required for nano-TSV connections removes lateral heat-spreading capability, leading to hotspot temperatures up to 45\% higher than equivalent front-side designs\cite{10925422}. Self-heating effects further complicate the picture at the 3\,nm node, where power-thermal co-design becomes essential to avoid reliability degradation\cite{10529350}. Additionally, the reduced DC resistance of BSPDNs shifts impedance resonance peaks into frequency ranges that may coincide with circuit switching activity\cite{8993617}, requiring careful decoupling capacitor placement during design. These challenges suggest that the electrical benefits of BSPDNs cannot be evaluated in isolation from their thermal and dynamic consequences.

This review has several limitations that should be acknowledged. The analysis is based on nine sources drawn exclusively from IEEE Xplore, which may not capture relevant work published in other databases or in non-English venues. Furthermore, the majority of the reviewed results are based on simulation rather than measured silicon data; only Intel's PowerVia study\cite{10185208} reports results from a fabricated test vehicle. The scope of the review is also limited to power integrity and routing metrics, and does not address manufacturing yield, cost, or supply chain considerations that will influence industrial adoption.

Future research should focus on three areas. First, measured thermal and electrical data from production-grade BSPDNs are needed to validate the simulation-based predictions that dominate the current literature. Second, power-thermal co-design methodologies that jointly optimise IR~drop and hotspot temperature at the block and system level remain an open challenge. Third, the extension of backside routing to include signal interconnects\cite{11192533} opens new possibilities for standard-cell scaling, but the associated design rules and parasitic trade-offs require further investigation. Addressing these questions will be critical to realising the full potential of backside power delivery in future technology nodes.


%TC:ignore

\section*{Acknowledgement}
\small
\textbf{Statement on the use of generative AI}\\
This article was written within the course EE4C01 Profile Orientation \& Academic Writing. Total word count: 2000.

When writing this text, we:
\begin{itemize}
    \item used generative AI to improve the grammar, style and/or spelling of the text.
\end{itemize}
\normalsize

\bibliography{bibliography}
%TC:endignore

\end{document}
